\documentclass[10pt,a4paper]{amsart}
\usepackage[margin=19mm]{geometry}
\usepackage{amsmath,amssymb,mathtools}
\usepackage[scr=rsfs,cal=euler]{mathalfa}
\usepackage{xfrac}
\usepackage[unicode,hidelinks,bookmarks=false]{hyperref}
\newcommand{\ie}{i.e.\ }
\newcommand{\cf}{cf.\ }
\newcommand{\resp}{respectively\ }
\newcommand{\Subsection}{\subsection}
\providecommand{\nobreakdash}{\nobreak\hbox{-}}
\setlength{\emergencystretch}{2em}
\multlinegap0pt
\usepackage{bm}
\usepackage{bbm}
\usepackage{dsfont}
\usepackage{xspace}
\usepackage{cancel}
\usepackage{mathtools}
\usepackage[shortlabels]{enumitem}
\usepackage{amsthm}
\makeatletter
\@ifundefined{theorem}{\newtheorem{theorem}{Theorem}}{}
\@ifundefined{lemma}{\newtheorem{lemma}[theorem]{Lemma}}{}
\@ifundefined{proposition}{\newtheorem{proposition}[theorem]{Proposition}}{}
\makeatother
\newcommand{\N}{\mathbb N}% by default, $\N$ is American naturals $\{1,2,\ldots\}$
\newcommand{\R}{\mathbb R}
\newcommand{\cN}{{\mathcal N}}
\title[Approximation by uniformly distributed sequences]{Approximation by uniformly distributed sequences}
\author{Gerardo Gonz\'alez Robert, Mumtaz Hussain, Nikita Shulga,, Benjamin Ward}
\date{Source: arXiv:2507.06583v1 (CC BY 4.0)}
\begin{document}
\maketitle

\noindent\emph{Source and licence.} Approximation by uniformly distributed sequences. Version arXiv:2507.06583v1 (CC BY 4.0), distributed under the Creative Commons Attribution 4.0 International licence, as recorded in the arXiv licence metadata for this exact version and shown on the abstract page https://arxiv.org/abs/2507.06583v1. Full rights notice: licenses/arxiv-2507.06583-CC-BY-4.0.txt.\par\smallskip

\noindent\textit{Editorial scope.} Counted unit: the Lemma whose source label is \texttt{uniform dirichlet} in the article (source0.tex lines 776--784, source label \texttt{uniform dirichlet}), delivered together with the article's own complete original proof of it (source0.tex lines 787--839, one \texttt{proof} environment of 53 lines). No mathematics is written, repaired or re-derived: statement and proof are the article's own text line for line. Required context retained uncounted: ubiquity statement (lines 734--744); Prop:AlmostCover (lines 751--761). Every definition, display and label of those lines is retained unchanged. Omitted from this package and not redistributed: 1--733, 745--750, 762--775, 785--786, 840--1355. Nothing inside a retained excerpt is omitted or altered: every retained body is delivered exactly as the article's lines are written, so no omission receipt of any kind accompanies this package. Citations: the retained excerpts carry no citation command at all, so no bibliography entry of the article is transcribed and no reference is redistributed. Reference adaptation: none was needed and none was made; every reference inside the delivered unit resolves inside it, so no unresolved reference or citation is produced. Printed slips kept literal: no printed word, symbol, hypothesis or number of the counted statement or of its proof was altered. Layout and class: the article's own class is \texttt{amsart} with packages including geometry, graphics, amsmath, amssymb, enumerate, comment, url, mathtools, graphicx, epsfig, hyperref, color, xcolor, pgf; the wrapper is the lane's pinned \texttt{amsart} editorial wrapper at 10pt on A4, which restates the theorem-like environments the retained text needs and adds the article's own macro definitions that the retained text needs (\texttt{\textbackslash N}, \texttt{\textbackslash R}, \texttt{\textbackslash cN}), with extra wrapper packages (bm, bbm, dsfont, xspace, cancel, mathtools, enumitem). The delivered numbering is the unit's own. Assets: no retained span contains a figure environment, a graphics-inclusion command, a PDF-page inclusion, a table or a TikZ picture; no image file is copied into this package, the package declares no asset dependency and no image file is redistributed with it. Licence: version arXiv:2507.06583v1 of this article is distributed under the Creative Commons Attribution 4.0 International licence, as recorded in the arXiv licence metadata for this exact version and shown on the abstract page https://arxiv.org/abs/2507.06583v1. A full rights notice is delivered at licenses/arxiv-2507.06583-CC-BY-4.0.txt. Characters: every retained excerpt is pure ASCII, so the delivered engine, XeLaTeX with its default Latin Modern text font, reports no missing character.
\begin{proposition} \label{ubiquity statement}
    Let $\omega$ be a $(\cN,v)$-d.s.s. . Let $\rho=(\rho_{1},\dots,\rho_{n})$ be an $n$-tuple of functions $\rho_{i}:\cN \to \R_{+}$ such that
    \begin{enumerate}
        \item[i)] each $\rho_{i}(N)\to 0$ as $N\to \infty$ ,
        \item[ii)] $\prod_{i=1}^{n}\rho_{i}(N)=v(N)$ for all $N\in\cN$ .
    \end{enumerate}
    Then there exists constant $c>0$ such that, for any ball $B\subset [0,1]^{n}$, there exists sufficiently large $k_{0}\in\N$ such that for all $k\geq k_{0}$
    \begin{equation*}
        \lambda_n\left( B\cap \bigcup_{N_{k-1}<j\leq N_{k}}\Delta\left(\boldsymbol{\omega}_{j},\rho(N_{k})\right)\right) \geq c\lambda_n(B).
    \end{equation*}
\end{proposition}

\begin{proposition}\label{Prop:AlmostCover}
Let $R\subseteq \mathbb{R}^n$ be a rectangle and $(\mathbf{r}_j)_{j\geq 1}$ a sequence in $\mathbb{R}_{+}^n$ converging to $(0,\ldots,0)$. For every $\varepsilon>0$ there exists $K=K(\varepsilon)\in\N$ such that for all $k\in \N_{\geq K}$ there is a finite set $\mathcal{C}=\mathcal{C}(k)\subseteq R$ verifying
\begin{enumerate}[i.]
\item If $\mathbf{y},\mathbf{w}\in\mathcal{C}$ are different, then $\Delta(\mathbf{y};\mathbf{r}_k)\cap\Delta(\mathbf{w};\mathbf{r}_k)=\varnothing$.
\item For each $\mathbf{y}\in\mathcal{C}$, the rectangle $\Delta(\mathbf{y},\mathbf{r}_k)$ is contained in the interior of $R$.
\item We have
\[
\lambda_n\left( R\setminus \bigcup_{\mathbf{y}\in\mathcal{C}} \Delta(\mathbf{y};\mathbf{r}_{k})\right)< \varepsilon.
\]
\end{enumerate}
\end{proposition}

\begin{lemma}\label{uniform dirichlet}
Let $\rho$ and $\omega$ satisfy the conditions stated in Proposition~\ref{ubiquity statement}. For any ball $B\subseteq [0,1]^{n}$ we have
\[
\lim_{k\to\infty}
\lambda_n\left( B\cap\bigcup_{j=1}^{N_k} \Delta(\boldsymbol{\omega}_j;\rho(N_k))\right)
=
\lambda_n(B).
\]
\end{lemma}

\begin{proof}
Take any $\varepsilon\in (0,1)$. Let $K$ be the natural number obtained from Proposition \ref{Prop:AlmostCover} applied on $R=B$ and $\mathbf{r}_k =\frac{1}{2} \rho(N_k)$. Take $k\in\N_{\geq K}$ and let $\mathcal{C}=\mathcal{C}(k)$ be the corresponding finite set. For every $\mathbf{y}\in \mathcal{C}$, we have 
\[
\left| 
\frac{A_{N_k}\left(\omega;\Delta\left(\mathbf{y}; \frac{\rho(N_k)}{2}\right) \right)}{N_k} -
 \lambda_n\left(\Delta\left(\mathbf{y}; \frac{\rho(N_k)}{2} \right)\right)
\right|
\leq 
D_{N_k} < v(N_k).
\]
Hence, since
\[
 \lambda_n\left(\Delta\left(\mathbf{y}; \frac{\rho(N_k)}{2} \right)\right) = v(N_k),
\]
we have
\[
A_{N_k}\left(\omega;\Delta\left(\mathbf{y}; \frac{\rho(N_k)}{2}\right) \right)\geq 1\, .
\]
So, for each $\mathbf{y}\in \mathcal{C}$, there exists 
\[
\omega(\mathbf{y})
\in 
\{\boldsymbol{\omega}_1,\ldots, \boldsymbol{\omega}_{N_k}\}\cap \Delta\left(\mathbf{y}; \frac{\rho(N_k)}{2} \right)\, .
\]
Fix one such point $\omega(\mathbf{y})$ for each $\mathbf{y}\in\mathcal{C}$. Note that 
\[
\Delta\left(\mathbf{y}; \frac{\rho(N_k)}{2} \right)
\subseteq
\Delta\left(\omega(\mathbf{y}); \rho(N_k) \right).
\]
Therefore, 
\begin{align*}
\lambda_n\left( B\cap\bigcup_{j=1}^{N_k} \Delta(\boldsymbol{\omega}_j;\rho(N_k))\right)
&\geq
\lambda_n\left( B\cap\bigcup_{\mathbf{y}\in \mathcal{C}} \Delta(x(\mathbf{y}) ;\rho(N_k))\right) \\
&\geq
\lambda_n\left( B\cap\bigcup_{\mathbf{y}\in \mathcal{C}} \Delta\left(\mathbf{y} ;\frac{\rho(N_k)}{2} \right)\right) \\
&>
(1-\varepsilon) \lambda_n(B).
\end{align*}
Since $\varepsilon>0$ was arbitrary, we conclude
\[
\liminf_{k\to\infty} \lambda_n\left( B\cap\bigcup_{j=1}^{N_k} \Delta(\boldsymbol{\omega}_j;\rho(N_k))\right)
\geq 
\lambda_n(B).
\]
The result follows, because we have for every $k\in\mathbb{N}$ that
\[
\lambda_n\left( B\cap\bigcup_{j=1}^{N_k} \Delta(\boldsymbol{\omega}_j;\rho(N_k))\right)
\leq 
\lambda_n(B).
\]
\end{proof}

\end{document}
